\section*{Chapitre \ref{ch:b3:membrane-traffic} --- Trafic membranaire et adressage des protéines}

\begin{solution}{exo:b3:membrane-traffic:1}
Tronçon hydrophobe amino-terminal~: \omterm{def:b3:membrane-traffic:signals}{peptide signal}, vers le réticulum
endoplasmique (et par défaut jusqu'à la surface). PKKKRKV~: \omterm{def:b3:membrane-traffic:signals}{signal de
localisation nucléaire}, vers le noyau par le pore. Hélice amphipathique
amino-terminale~: \omterm{def:b3:membrane-traffic:signals}{préséquence} mitochondriale, vers la matrice
mitochondriale. KDEL~: récupération depuis le Golgi vers le réticulum.
Mannose-6-phosphate~: du Golgi trans vers l'\omterm{def:b3:membrane-traffic:endocytosis}{endosome} et le \omterm{def:b3:membrane-traffic:endocytosis}{lysosome}.
\end{solution}

\begin{solution}{exo:b3:membrane-traffic:2}
(1) Sans membranes, la chaîne traduite était plus longue d'une
vingtaine de résidus que la protéine sécrétée~: un segment est retiré
au cours de la sécrétion. (2) Avec des membranes présentes pendant la
traduction, le produit avait la taille mature et était protégé de la
protéase~: il était entré dans les vésicules et y avait perdu le
segment. (3) Des membranes ajoutées après la traduction ne faisaient ni
l'un ni l'autre~: l'entrée est couplée à la synthèse ---
cotraductionnelle --- et le segment, le \omterm{def:b3:membrane-traffic:signals}{peptide signal}, n'agit que sur
une chaîne naissante.
\end{solution}

\begin{solution}{exo:b3:membrane-traffic:3}
Réticulum vers Golgi~: \omterm{def:b3:membrane-traffic:coats}{COPII}, vers l'avant. Golgi vers réticulum~:
\omterm{def:b3:membrane-traffic:coats}{COPI}, rétrograde. Golgi trans vers \omterm{def:b3:membrane-traffic:endocytosis}{endosome}~: \omterm{def:b3:membrane-traffic:coats}{clathrine} et ses
adaptateurs, vers l'avant, en direction du \omterm{def:b3:membrane-traffic:endocytosis}{lysosome}. Membrane plasmique
vers \omterm{def:b3:membrane-traffic:endocytosis}{endosome}~: \omterm{def:b3:membrane-traffic:coats}{clathrine}, vers l'intérieur (endocytose).
\end{solution}

\begin{solution}{exo:b3:membrane-traffic:4}
Les sucres ne sont ajoutés que dans la lumière du réticulum et du
Golgi. La face luminale de chaque vésicule reste luminale au
bourgeonnement comme à la fusion, et lorsqu'une vésicule fusionne avec
la membrane plasmique sa face luminale devient l'extérieur de la
cellule~; l'orientation n'est jamais inversée, de sorte que ce qui a été
glycosylé dedans se retrouve dehors.
\end{solution}

\begin{solution}{exo:b3:membrane-traffic:5}
$t^{*} = \ln(0.05/0.2)/(0.05 - 0.2) = \ln 4/0.15 = \qty{9.2}{min}$~;
$B(t^{*}) = (0.2/(-0.15))(e^{-1.85} - e^{-0.46}) = 0.63$. À \qty{60}{min},
$A \approx 0$, $B = 1.33\,(e^{-3} - e^{-12}) = 0.066$, donc $C = 0.93$.
\end{solution}

\begin{solution}{exo:b3:membrane-traffic:6}
Le \omterm{def:b3:membrane-traffic:signals}{peptide signal} met l'extrémité aminée dans la lumière~; chaque
tronçon hydrophobe est tour à tour un signal d'arrêt et un signal de
reprise de transfert, de sorte qu'il y a quatre hélices
transmembranaires (environ les résidus 30--52, 90--112, 150--172,
210--232). Segments~: extrémité aminée luminale~; 52--90 cytosolique~;
112--150 luminale~; 172--210 cytosolique~; extrémité carboxylée
luminale. Glycosylation sur les seuls segments luminaux ---
l'extrémité aminée, 112--150 et l'extrémité carboxylée --- qui
deviennent extracellulaires à la surface.
\end{solution}

\begin{solution}{exo:b3:membrane-traffic:7}
$J_{\max} = \num{30000}\times 0.25\times 0.05/0.30 = 1250$ par minute~;
fraction de surface $k_{r}/(k_{i} + k_{r}) = 0.05/0.30 = 0.17$. Doubler
$R$ double exactement $J$~; c'est la vitesse de retour qui est le
goulot (la plupart des récepteurs sont à l'intérieur), et doubler
$k_{r}$ donnerait $\num{30000}\times 0.25
\times 0.1/0.35 \approx 2140$, un facteur $1.7$.
\end{solution}

\begin{solution}{exo:b3:membrane-traffic:8}
(a) Pas de récepteur~: les hydrolases marquées suivent la voie par
défaut et sont sécrétées~; les \omterm{def:b3:membrane-traffic:endocytosis}{lysosomes} manquent d'enzymes et se
remplissent de matériel non digéré --- une maladie de surcharge.
(b) Pas de phosphotransférase~: le même phénotype par une autre lésion
(maladie des cellules à inclusions), les enzymes non marquées et
sécrétées. (c) \omterm{def:b3:membrane-traffic:endocytosis}{Endosomes} neutralisés~: le récepteur ne peut pas lâcher
sa cargaison, il n'est donc pas recyclé et les hydrolases sont mal
livrées ou sécrétées~; la LDL et d'autres récepteurs cessent aussi de
tourner.
\end{solution}

\begin{solution}{exo:b3:membrane-traffic:9}
La tyrosine appartient au motif de la queue que reconnaît l'adaptateur
de la \omterm{def:b3:membrane-traffic:coats}{clathrine}~; sans elle le récepteur n'est pas rassemblé dans les
puits à \omterm{def:b3:membrane-traffic:coats}{manteau}. Les récepteurs fixent normalement la LDL mais sont
répartis uniformément sur la surface, exclus des puits, et le ligand
reste dehors --- une maladie de l'adressage, non de la fixation.
\end{solution}

\begin{solution}{exo:b3:membrane-traffic:10}
Avec $k_{1} = k_{2} = k$, essayons $B = kt\,e^{-kt}$~:
$\mathrm{d}B/\mathrm{d}t
= k e^{-kt} - k^{2}t e^{-kt} = kA - kB$, comme il se doit, avec
$B(0) = 0$. Maximum là où $1 - kt = 0$~: $t^{*} = 1/k$,
$B = e^{-1} = 0.37$. En général
$C(t) = 1 - (k_{2}e^{-k_{1}t} - k_{1}e^{-k_{2}t})/(k_{2} - k_{1})$, qui
est symétrique en échangeant $k_{1}$ et $k_{2}$~: la seule courbe
granulaire ne peut pas dire lequel des deux compartiments est le
rapide.
\end{solution}

\begin{solution}{exo:b3:membrane-traffic:11}
Chaque vésicule a pour aire $4\pi(50\,\text{nm})^{2} = \qty{0.031}{\micro
m^2}$~; un millier par minute fait \qty{31}{\micro m^2/min}, donc toute
la surface en $3000/31 \approx \qty{95}{min}$. La cellule ne rétrécit
pas parce que la membrane revient par exocytose (vésicules de
recyclage) à la même vitesse --- un équilibre de flux. Chaque vésicule
contient $5\times 10^{-19}$ L~; un millier par minute pendant une heure
fait $3\times 10^{-14}$ L, soit environ \qty{1.5}{\%} du volume d'une
cellule de \qty{2}{pL} par heure.
\end{solution}

\begin{solution}{exo:b3:membrane-traffic:12}
La toxine botulique agit dans la terminaison motrice~: aucune
acétylcholine n'est libérée, le muscle ne reçoit aucun ordre et reste
mou --- paralysie flasque. La toxine tétanique est remontée le long de
l'axone moteur puis passe dans les interneurones inhibiteurs qui
retiennent normalement les motoneurones~; leur \omterm{def:b3:membrane-traffic:fusion}{SNARE} clivée, ni glycine
ni GABA n'est libéré, les motoneurones déchargent sans frein et tous
les muscles se contractent --- paralysie spastique. Quelques
nanogrammes de toxine botulique injectés dans un muscle trop actif le
font taire localement pendant des mois (jusqu'à ce que de la \omterm{def:b3:membrane-traffic:fusion}{SNARE}
neuve soit produite et que la terminaison récupère)~: un traitement des
dystonies, de la spasticité, du strabisme et des rides.
\end{solution}

\begin{solution}{pb:b3:membrane-traffic:1}
\textbf{1.} $10^{7}\times \qty{25000}{Da}\times 1.66\times 10^{-24}$ g
$= \qty{0.42}{pg}$ par minute, \qty{600}{pg} par jour --- le double de
son propre contenu en protéines.
\textbf{2.} Volume d'une vésicule $\tfrac{4}{3}\pi\,35^{3} = 1.8\times 10^{5}$
nm$^{3}$, dont \qty{30}{\%} font $5.4\times 10^{4}$~; une molécule
d'enzyme $\tfrac{4}{3}\pi\,2^{3} = 34$ nm$^{3}$~: environ \num{1600}
molécules par vésicule~; $10^{7}/1600 \approx 6200$ vésicules par
minute.
\textbf{3.} Aire $4\pi\,35^{2} = \qty{0.0154}{\micro m^2}$ chacune, soit \qty{95}{\micro
m^2} par minute~: les \qty{6000}{\micro m^2} du réticulum en une heure
environ si rien ne revenait.
\textbf{4.} $\qty{95}{\micro m^2} = 9.5\times 10^{7}$ nm$^{2}$, deux feuillets
à \qty{0.7}{nm^2}~: $2.7\times 10^{8}$ lipides par minute. La cellule
ramène de la membrane par les vésicules \omterm{def:b3:membrane-traffic:coats}{COPI} et le recyclage, et
synthétise du lipide pour remplacer ce qui part définitivement dans les
granules.
\textbf{5.} Volume d'un granule $5.2\times 10^{8}$ nm$^{3}$, dont \qty{30}{\%} font
$1.6\times 10^{8}$~: $4.7\times 10^{6}$ molécules par granule~; $6\times
10^{8}$ molécules par heure remplissent environ $130$ granules.
\textbf{6.} Chaque granule ajoute $\pi d^{2} = \qty{3.1}{\micro m^2}$~; $300$
ajoutent \qty{940}{\micro m^2} à \qty{50}{\micro m^2}~: un facteur
vingt. La membrane doit être récupérée par une endocytose compensatrice
en quelques minutes.
\textbf{7.} $t^{*} = \ln 2/0.05 = \qty{13.9}{min}$, $B = 0.50$.
\textbf{8.} $A = e^{-3} = 0.05$~; $B = 2(e^{-1.5} - e^{-3}) = 0.35$~;
$C = 0.60$.
\textbf{9.} \qty{10}{min} et \qty{20}{min}~; total moyen \qty{30}{min}.
\textbf{10.} $t^{*} = \ln(0.1)/(-0.09) = \qty{25.6}{min}$~; $B(t^{*}) =
(0.1/(-0.09))(e^{-2.56} - e^{-0.256}) = 0.77$~: le Golgi enfle, gorgé
des trois quarts de la cargaison.
\textbf{11.} Au maximum $k_{1}e^{-k_{1}t^{*}} = k_{2}e^{-k_{2}t^{*}}$~;
en reportant dans $B$ on obtient $B_{\max} = (k_{1}/k_{2})^{\,k_{2}/(k_{2} -
k_{1})}$~: la puissance est $k_{2}/(k_{2} - k_{1})$. Pour
$k_{1} > k_{2}$ l'exposant est négatif et la base supérieure à $1$~;
pour $k_{1} < k_{2}$ la base est inférieure à $1$ et l'exposant positif
--- dans les deux cas la valeur est inférieure à $1$ (vérification~:
$2^{-1} = 0.5$ ici).
\textbf{12.} Pour $k_{1} \gg k_{2}$, $B(t) \to e^{-k_{2}t}$~: le
marqueur est dans le Golgi d'emblée et en sort à $k_{2}$ --- l'étape du
réticulum est invisible et la courbe golgienne est une simple
décroissance.
\textbf{13.} $\theta = 2/(2+2) = 0.5$~; $J = \num{50000}\times 0.2\times
0.1\times 0.5/(0.1 + 0.1) = 2500$ particules par minute, \num{150000}
par heure.
\textbf{14.} $1.5\times 10^{5}\times 1500 = 2.3\times 10^{8}$ molécules
de cholestérol par heure~; $5\times 10^{9}/2.3\times 10^{8} \approx 22$
heures.
\textbf{15.} Fraction de surface $k_{r}/(k_{r} + k_{i}\theta) = 0.5$~;
un aller-retour prend $1/(k_{i}\theta) + 1/k_{r} = 10 + 10 = \qty{20}{min}$~:
environ $60$ voyages en vingt heures.
\textbf{16.} $J = 1250$ par minute. En deçà de la saturation
$J \propto R\,[L]$, de sorte qu'avec la moitié des récepteurs la LDL
plasmatique doit doubler pour la même épuration.
\textbf{17.} $R$ revenu à \num{50000}~: l'épuration est rétablie, la LDL
plasmatique retombe vers la normale --- l'effet des statines.
\textbf{18.} Sans récepteurs, la LDL n'est épurée que par des voies non
spécifiques et lentes, de sorte qu'elle circule des jours au lieu
d'heures, s'oxyde, et est captée par les récepteurs éboueurs des
macrophages, qui deviennent les cellules spumeuses des plaques
d'athérome~; et les statines, qui agissent en induisant des récepteurs,
n'ont rien à induire.
\textbf{19.} $V = \tfrac{4}{3}\pi(0.25)^{3} = \qty{0.065}{\micro m^3} =
6.5\times 10^{-17}$ L. À pH $7.2$~: $6.3\times 10^{-8}\times 6.5\times
10^{-17} = 4\times 10^{-24}$ mol --- environ $2$ protons libres. À pH
$4.7$~: $2\times 10^{-5}\times 6.5\times 10^{-17} = 1.3\times 10^{-21}$
mol, soit environ $780$.
\textbf{20.} $50\times 780 \approx \num{39000}$ protons~; $50$ pompes à
$200$ par seconde en déplacent \num{10000} par seconde~: environ
\qty{4}{s}.
\textbf{21.} Faire entrer de la charge positive rend la lumière
positive, et quelques milliers de charges arrêteraient la pompe~; du
chlorure entre par un canal (ou des cations sortent) pour neutraliser
la charge.
\textbf{22.} Des enzymes inactives à pH neutre ne font aucun mal si un
\omterm{def:b3:membrane-traffic:endocytosis}{lysosome} fuit ou si elles sont mal adressées au cytosol ou à la
surface~: l'exigence d'acidité confine la digestion au compartiment
bâti pour elle.
\textbf{23.} Les deux doivent porter un capteur de pH --- des
histidines, dont la protonation vers pH~6 change la surface de
fixation~: le récepteur de la LDL rabat un domaine sur son propre site
de fixation une fois protoné, le récepteur du mannose-6-phosphate perd
son affinité de la même façon.
\textbf{24.} $10^{7.2 - 4.7} = 10^{2.5} \approx 300$ fois plus
concentrée~; la base piégée consomme des protons, élève le pH
lysosomal et inactive les hydrolases --- c'est ainsi que la chloroquine
tue le parasite du paludisme dans sa vacuole digestive et qu'elle bloque
l'\omterm{def:b3:membrane-traffic:endocytosis}{autophagie} dans les expériences.
\textbf{25.} Maximum golgien à \qty{14}{min}~; quelque \num{6200}
vésicules quittent le réticulum par minute~; \num{150000} particules de
LDL captées par heure.
\end{solution}
